\begin{table}[htpb]
    \centering
    \caption{Degrees of freedom across lag orders}
    \label{tab:df_budget}
    \begin{tabular}{lccc}
        \toprule
         & \textbf{VAR $p=4$} & \textbf{VAR $p=8$} & \textbf{VAR $p=12$} \\
        \midrule
        Usable observations $T_u$ & 167 & 163 & 159 \\
        Coefficients per equation $k$ & 41 & 81 & 121 \\
        Degrees of freedom $T_u-k$ & 126 & 82 & 38 \\
        Observations per parameter $T_u/k$ & 4.07 & 2.01 & 1.31 \\
        Coefficients, full system & 410 & 810 & 1{,}210 \\
        \bottomrule
    \end{tabular}

    \vspace{0.15cm}
    \raggedright
    \footnotesize{\textit{Note:} $T$ denotes the raw sample size and $T_u$ the usable sample size, $T_u = T-p$; $k$ denotes the number of regressors per equation, $k = 10p+1$ ($K=10$ variables, $p$ lags, one constant). The raw sample contains $T=171$ quarterly observations (1959:Q2--2001:Q4), of which $p$ are consumed as initial conditions, so that $T_u = T-p$. Ordinary least squares requires $k < T_u$, that is $10p+1 < 171-p$, that is $p \leq 15$; at $p=15$ only 5 degrees of freedom per equation remain.}
\end{table}
